\chapter{The $\kappa$-conductor scalar packages}
\label{ch:kappa-conductor}

Four constructions meet at the scalar end of the Open Beilinson tower:
\[
 K^c(A,A^!)=c(A)+c(A^!),\qquad
 K^\kappa(A,A^!)=\kappa(A)+\kappa(A^!),
\]
\[
 K_{\mathrm{ghost}}^{\mathrm{leg}}(\mathcal P)
 =-c_{\mathrm{ghost}}(\mathcal P),\qquad
 K_E(A)=\chi\bigl(\bar B_X(A)\bigr).
\]
Their type signatures determine their meaning.  The first receives a
pair of conformal vertex algebras.  The second also receives a chosen
modular-characteristic normalization.  The third receives a concrete
conformal BRST presentation~$\mathcal P$.  The fourth receives a
completed bar object together with a convergence theorem.  The
universal conductor of Chapter~\ref{chap:universal-conductor} is the
ordered-to-symmetric averaging morphism; these four numbers are scalar
shadows available after further choices.

\section{The first calculation: the Virasoro pair}

The Virasoro family displays the separation in its smallest useful
form.  For the companion central charges $c$ and $26-c$,
\[
 K^c=c+(26-c)=26,
 \qquad
 K^\kappa=\frac c2+\frac{26-c}{2}=13.
\]
The standard reparametrisation presentation has one fermionic
$bc(2)$ pair.  Since
\[
 c_{bc}(\lambda)=-2(6\lambda^2-6\lambda+1),
\]
its ghost scalar is
\[
 K_{\mathrm{ghost}}^{\mathrm{leg}}
 =-c_{bc}(2)=2(24-12+1)=26.
\]
Thus the Virasoro gauge presentation yields the ordered triple
\[
 (K^c,K^\kappa,K_{\mathrm{ghost}}^{\mathrm{leg}})=(26,13,26).
\]
The equality of the first and third entries is the direct anomaly
calculation for this presentation; the middle entry carries the
modular-characteristic normalization $\kappa(\Vir_c)=c/2$.

The normalizations used throughout the chapter are those of
Table~\ref{tab:master-invariants}:
\[
\kappa(\mathcal H_{\ell}^{\oplus d})=d\ell,
\qquad
\kappa(V_k(\mathfrak g))
=\frac{\dim(\mathfrak g)(k+h^\vee)}{2h^\vee},
\]
\[
\kappa(\Vir_c)=\frac c2.
\]
Affine Kac--Moody and Heisenberg companion sums vanish in this scalar
normalization.  Under the genus-one diagonal-trace package
$H_{\mathrm{diag}}^{g=1}$,
\[
\kappa(\mathcal W_N)=c(H_N-1).
\]
Together with $H_{W_N}^{\mathrm{DS/bar}}$, this gives
$K^\kappa=(H_N-1)K^c$.  The K3/BKM row carries subscripted invariants
whose domains are specified in
Section~\ref{sec:conductor-cross-volume}.

\section{The Euler characteristic of the reduced bar object}
\label{sec:conductor-euler-pairing}

Let $A$ be an augmented complete chiral algebra on a smooth projective
curve~$X$.  Suppose every conformal-weight piece of
$H^\bullet(\bar B_X(A))$ is finite-dimensional and the alternating
weightwise sum converges in a specified completion.  Then
\[
 K_E(A):=\chi\bigl(\bar B_X(A)\bigr)
 =\sum_{n\in\mathbb Z}(-1)^n
 \dim H^n\bigl(\bar B_X(A)\bigr)
\]
is a bar-complex invariant.  An identity between $K_E$ and one of the
three conformal scalars requires a comparison map from the completed
bar complex to the relevant conformal or BRST complex.  The comparison
map and its convergence package are part of the theorem statement.

\section{Central-charge and modular-characteristic sums}
\label{sec:conductor-central-charge}

For a specified conformal companion pair $(A,A^!)$, set
\[
 K^c(A,A^!):=c(A)+c(A^!),\qquad
 K^\kappa(A,A^!):=\kappa(A)+\kappa(A^!).
\]
When both members have a common level-independent ratio
$\varrho=\kappa/c$, direct addition gives
\begin{equation}
 \label{eq:kappa-conductor-ratio-bridge}
 K^\kappa(A,A^!)=\varrho K^c(A,A^!).
\end{equation}
The formula is scalar algebra.  Its interpretation as Verdier--Koszul
complementarity additionally uses the named theorem identifying the
second member with the algebra-level companion.

For the Virasoro companion and the principal DS reflection one obtains
\[
K^c(\Vir_c,\Vir_{26-c})=26,
\qquad
c_N(k)+c_N(-k-2N)=4N^3-2N-2.
\]
The package $H_{W_N}^{\mathrm{DS/bar}}$ identifies the second reflected
central-charge function with the principal companion and promotes its
sum to $K^c(\mathcal W_N,\mathcal W_N^!)$.
The Bershadsky--Polyakov calculation fixes the central-charge lane and
isolates the remaining genus-one problem.  Fehily--Kawasetsu--Ridout
define an ordinary bosonic vertex algebra: the generators $J,T,G^+$,
and $G^-$ are all even, with conformal weights $1,2,3/2,3/2$.
They use the conformal vector
\cite[Eq.~\textup{(}2.2\textup{)}]{FKR20} for which
\begin{equation}
 \label{eq:kappa-conductor-bp-standard-c}
 c_{\mathrm{BP}}(k)
 =-\frac{(2k+3)(3k+1)}{k+3}.
\end{equation}
Writing $t=k+3$ gives
\[
 c_{\mathrm{BP}}(k)=25-6t-\frac{24}{t}.
\]
The involution $k\mapsto-k-6$ sends $t$ to $-t$, hence
\begin{equation}
 \label{eq:kappa-conductor-bp-standard-sum}
 c_{\mathrm{BP}}(k)+c_{\mathrm{BP}}(-k-6)=50
 \qquad (k\ne-3).
\end{equation}
Thus $K^c_{\mathrm{BP}}=50$ is an exact rational-function identity.
The modular characteristic $\kappa_{\mathrm{BP}}$ is
\ClaimStatusOpen.  Its construction requires the genus-one
nonseparating curvature of the minimal Drinfeld--Sokolov complex,
including the charged ghosts, neutral fields, conformal improvement,
and mixed channels.  The source-correct reciprocal-weight diagnostic is
\[
 1+\frac23+\frac23+\frac12=\frac{17}{6};
\]
it records the four even strong generators and supplies a check on that
future calculation.  A completed calculation establishing
$\kappa_{\mathrm{BP}}(k)=c_{\mathrm{BP}}(k)/6$ would imply
\begin{equation}
 \label{eq:kappa-conductor-bp-standard-kappa}
 K^\kappa_{\mathrm{BP}}=\frac16K^c_{\mathrm{BP}}=\frac{25}{3}.
\end{equation}
Equation~\eqref{eq:kappa-conductor-bp-standard-kappa} is a conditional
implication recorded in the open conductor ledger.

The algebra-level assertion
$\mathrm{BP}_k^!\simeq\mathrm{BP}_{-k-6}$ carries the package
\[
 H_{\mathrm{BP}}^{\mathrm{DS/bar}}:
 \quad k\ne-3,\quad
 \text{finite-type bar dual or continuous completed Verdier pairing},
 \quad\text{subregular DS/bar transport}.
\]
Equation~\eqref{eq:kappa-conductor-bp-standard-sum} is an identity of
rational functions; $H_{\mathrm{BP}}^{\mathrm{DS/bar}}$ promotes its
two terms to a same-family Verdier--Koszul pair.

The secondary rational function recorded elsewhere in the manuscript is
\[
 c_{\mathrm{BP}}^{\mathrm{shift}}(k)
 =2-\frac{24(k+1)^2}{k+3}.
\]
It satisfies
\begin{equation}
 \label{eq:kappa-conductor-bp-shifted-sum}
 c_{\mathrm{BP}}^{\mathrm{shift}}(k)
 +c_{\mathrm{BP}}^{\mathrm{shift}}(-k-6)=196.
\end{equation}
This is the computed-secondary shifted scalar lane.  The standard
Fehily--Kawasetsu--Ridout conformal vector supplies
Equations~\eqref{eq:kappa-conductor-bp-standard-c} and
\eqref{eq:kappa-conductor-bp-standard-sum}; the modular equation
\eqref{eq:kappa-conductor-bp-standard-kappa} remains conditional on its
genus-one curvature calculation.

\section{The scalar attached to a conformal ghost presentation}
\label{sec:conductor-ghost-charge}

Fix a finite conformal ghost presentation $\mathcal P$.  A fermionic
$bc(\lambda)$ pair contributes
\[
 c_{bc}(\lambda)=-2(6\lambda^2-6\lambda+1),
\]
and a bosonic $\beta\gamma(\lambda)$ pair contributes the opposite
central charge.  If $s_\alpha=+1$ for a fermionic $bc$ pair and
$s_\alpha=-1$ for a bosonic $\beta\gamma$ pair, define
\begin{equation}
 \label{eq:kappa-conductor-ghost-definition}
 K_{\mathrm{ghost}}^{\mathrm{leg}}(\mathcal P)
 :=-c_{\mathrm{ghost}}(\mathcal P)
 =\sum_\alpha s_\alpha
 2(6\lambda_\alpha^2-6\lambda_\alpha+1).
\end{equation}
The input~$\mathcal P$ includes the conformal vector, the charged and
neutral fields, their parities, and the improvement term.  A change of
presentation can change the scalar.  The level-$1$
Frenkel--Kac--Segal equivalence illustrates this dependence: a
Wakimoto gauge presentation counts gauge pairs, while the lattice-VOA
presentation counts matter bosons.

For minimal Drinfeld--Sokolov reduction of
$V_k(\mathfrak{sl}_3)$, the standard affine central charge is
$c_{\mathrm{aff}}(k)=8k/(k+3)$.  Equation~\eqref{eq:kappa-conductor-bp-standard-c}
therefore gives the exact total conformal reduction shift
\begin{equation}
 \label{eq:kappa-conductor-bp-total-shift}
 \Delta^{\mathrm{DS}}_{\mathrm{BP}}(k)
 :=c_{\mathrm{BP}}(k)-c_{\mathrm{aff}}(k)
 =-(6k+1).
\end{equation}
Its decomposition into charged ghosts, neutral fields, and conformal
improvement is presentation data.  Accordingly,
Equation~\eqref{eq:kappa-conductor-bp-total-shift}, the canonical
conductor~$50$, and the shifted secondary sum~$196$ occupy three
separate scalar lanes.

\section{The comparison theorem}
\label{sec:trinity-theorem}

\begin{theorem}[Typed scalar-conductor comparison]
\label{thm:conductor-trinity}
\ClaimStatusConditional
The type signature is
\[
 (\text{Open quadrant},\ \text{conformal/BRST scalar presentation},\
 \text{level }5,\ H_{\mathrm{pair}}\cup H_{\mathrm{BRST}}).
\]
Let $(A,A^!)$ be a conformal companion pair with a common anomaly ratio
$\varrho$, and let $\mathcal P$ be a specified finite conformal ghost
presentation.  Then:
\begin{enumerate}
\item the scalar companion identity is
$K^\kappa(A,A^!)=\varrho K^c(A,A^!)$;
\item an equality
$K_{\mathrm{ghost}}^{\mathrm{leg}}(\mathcal P)=K^c(A,A^!)$
is proved by the direct spin sum~\eqref{eq:kappa-conductor-ghost-definition}
in the chosen presentation;
\item the standard Virasoro and principal $\mathcal W_N$ gauge
presentations satisfy this equality;
\item the rank-one abelian-affine presentation has
$(K^c,K^\kappa,K_{\mathrm{ghost}}^{\mathrm{leg}})=(0,0,2)$;
\item on $H_{\mathrm{BP}}^{\mathrm{DS/bar}}$, the standard
Bershadsky--Polyakov central-charge pair has $K^c=50$, while
$\Delta^{\mathrm{DS}}_{\mathrm{BP}}(k)=-(6k+1)$ records the total
conformal reduction shift; its modular conductor $K^\kappa$ is the open
genus-one quantity described above.
\end{enumerate}
\end{theorem}

\begin{proof}
The first clause is Equation~\eqref{eq:kappa-conductor-ratio-bridge}.
For Virasoro, the $bc(2)$ calculation gives~$26$.  For principal
$\mathcal W_N$, one fermionic $bc(j)$ pair for each spin
$j=2,\ldots,N$ gives
\[
 \sum_{j=2}^{N}2(6j^2-6j+1)=4N^3-2N-2,
\]
which is the central-charge companion sum on that witness lane.  A
rank-one $bc(1)$ pair gives~$2$, while the Heisenberg scalar companion
sums vanish in the curved convention of
Theorem~\ref{thm:scalar-trace-complementarity}.  The BP
clauses are Equations~\eqref{eq:kappa-conductor-bp-standard-sum}
and~\eqref{eq:kappa-conductor-bp-total-shift}, with the categorical
interpretation supplied by $H_{\mathrm{BP}}^{\mathrm{DS/bar}}$.  The
modular clause becomes numerical after the genus-one curvature
calculation defining $\kappa_{\mathrm{BP}}$.
\end{proof}

\section{Finite ghost spin sums}
\label{sec:platonic-conductor}

\begin{theorem}[Finite $bc$-spin sum]
\label{thm:platonic-conductor}
\ClaimStatusProvedHere
The type signature is
\[
 (\text{Open quadrant},\ \text{finite conformal ghost presentation},\
 \text{level }5,\ H_{\mathrm{BRST}}^{\mathrm{fin}}).
\]
For a finite presentation consisting of fermionic $bc$ pairs of
weights $\lambda_\alpha$,
\[
 K_{\mathrm{ghost}}^{\mathrm{leg}}(\mathcal P)
 =\sum_\alpha2(6\lambda_\alpha^2-6\lambda_\alpha+1).
\]
Consequently:
\begin{enumerate}
\item a rank-one $bc(1)$ presentation has scalar~$2$;
\item a presentation with $d$ copies of $bc(1)$ has scalar~$2d$;
\item the Virasoro $bc(2)$ presentation has scalar~$26$;
\item the principal $\mathcal W_N$ gravity presentation with spins
$2,\ldots,N$ has scalar $4N^3-2N-2$.
\end{enumerate}
\end{theorem}

\begin{proof}
Each clause follows by substitution in
Equation~\eqref{eq:kappa-conductor-ghost-definition}.  For the final
clause,
\begin{align*}
\sum_{j=2}^{N}2(6j^2-6j+1)
&=12\left(\frac{N(N+1)(2N+1)}6-1\right)\\
&\quad-12\left(\frac{N(N+1)}2-1\right)+2(N-1)\\
&=4N^3-2N-2.
\end{align*}
\end{proof}

\section{Per-family corollaries}
\label{sec:conductor-corollaries}

\begin{corollary}[Heisenberg]
\label{cor:K-heisenberg}
In the rank-one abelian-affine gauge presentation,
\[
K_{\mathrm{ghost}}^{\mathrm{leg}}=2,
\qquad K^c=K^\kappa=0.
\]
The first entry is the $bc(1)$ spin sum.  The other two entries are the
central and modular-characteristic companion sums in the curved scalar
convention.
\end{corollary}

\begin{corollary}[Affine Kac--Moody]
\label{cor:K-affine-KM}
Let $\mathfrak g$ be simple and $k\ne-h^\vee$.  For the reflected
companion level $k'=-k-2h^\vee$,
\[
K^c(V_k(\mathfrak g),V_{k'}(\mathfrak g))=2\dim\mathfrak g,
\qquad K^\kappa=0.
\]
In the gauge presentation with one fermionic $bc(1)$ pair per Lie
direction,
\[
K_{\mathrm{ghost}}^{\mathrm{leg}}=2\dim\mathfrak g.
\]
The first identity follows from the Sugawara central charge; the second
follows because $k+h^\vee$ changes sign under the level reflection.
\end{corollary}

\begin{corollary}[Virasoro]
\label{cor:K-virasoro}
For the Virasoro companion pair,
\[
 K^c=26,\qquad K^\kappa=13,
 \qquad K_{\mathrm{ghost}}^{\mathrm{leg}}(bc(2))=26.
\]
The scalar midpoint is $c=13$, where each modular characteristic is
$13/2$.  The string anomaly-cancellation point uses matter central
charge~$26$ together with the reparametrisation ghost central
charge~$-26$.
\end{corollary}

\begin{corollary}[$W_N$ cubic; \ClaimStatusConditional]
\label{cor:K-WN-cubic}
Type signature: \textup{(}Open quadrant, principal DS and finite ghost
presentations, Beilinson level~$5$, hypothesis package
$H_{\mathrm{diag}}^{g=1}+H_{W_N}^{\mathrm{DS/bar}}$ for the modular
companion statement\textup{)}.
For principal $\mathcal W_N$ with the standard gravity ghosts at spins
$2,\ldots,N$, the exact identities are
\[
 K_{\mathrm{ghost}}^{\mathrm{leg}}
 =4N^3-2N-2,
\qquad
c_N(k)+c_N(-k-2N)=4N^3-2N-2.
\]
On $H_{W_N}^{\mathrm{DS/bar}}$, the second identity is the companion
statement $K^c=4N^3-2N-2$.  On
$H_{\mathrm{diag}}^{g=1}+H_{W_N}^{\mathrm{DS/bar}}$,
\[
 K^\kappa=(H_N-1)(4N^3-2N-2).
\]
The third forward differences are
\[
\Delta^3K^c=24,
\qquad
\Delta^3\left(\frac{K^c}{2}\right)=12.
\]
The scalar midpoint is
\[
c_*(\mathcal W_N)=\frac{K^c}{2}=2N^3-N-1.
\]
For $N=2,3,4,5,6$ these values are respectively
$13,50,123,244,425$.  For principal $\mathcal W_3$,
\[
c(k)=2-\frac{24(k+2)^2}{k+3}
=50-24\left((k+3)+(k+3)^{-1}\right).
\]
Hence $c(k)=50$ is equivalent to $(k+3)^2+1=0$, and the two
analytically continued midpoint levels are
\[
k=-3\pm i.
\]
\end{corollary}

\begin{corollary}[Bershadsky--Polyakov]
\label{cor:K-BP}
\ClaimStatusProvedHere
For $k\ne-3$, the standard Fehily--Kawasetsu--Ridout conformal vector
satisfies
\[
c_{\mathrm{BP}}(k)+c_{\mathrm{BP}}(-k-6)=50.
\]
On $H_{\mathrm{BP}}^{\mathrm{DS/bar}}$ this is the same-family
central-charge value $K^c=50$.  The secondary shifted rational function
has companion sum~$196$.  The total conformal shift of the minimal DS
reduction is $-(6k+1)$.

The modular characteristic is \ClaimStatusOpen.  If its genus-one
curvature calculation yields
$\kappa_{\mathrm{BP}}=c_{\mathrm{BP}}/6$, then
$K^\kappa_{\mathrm{BP}}=25/3$ follows from the proved central-charge
identity.  The even generator parities give the independent diagnostic
$17/6$ for the reciprocal-weight sum entering that calculation.

The standard midpoint $c=25$ occurs at
\[
k=-3\pm2i,
\]
and the involution $k\mapsto-k-6$ exchanges the two levels.
\end{corollary}

\begin{remark}[The integer $24$ in the third difference]
\label{rem:third-difference-24}
The identity
$\Delta^3(4N^3-2N-2)=24$ is the finite-difference signature of the
principal $W_N$ ghost spin sum.  Its content is the elementary rule
$\Delta^3N^3=6$.  The K3/BKM package obtains its integers from a
different construction:
\[
\kappa_{\mathrm{fiber}}(K3)=24,
\qquad
\kappa_{\mathrm{BKM}}(\Delta_5)=5,
\qquad
\kappa_{\mathrm{BKM}}(\Phi_{12})=12.
\]
The shared integer~$24$ is therefore an arithmetic meeting point of
two separately typed invariants.
\end{remark}

\section{The extension frontier}
\label{sec:conductor-exclusion-frontier}

Equation~\eqref{eq:kappa-conductor-ghost-definition} applies to a
specified finite conformal ghost presentation.  Three broader classes
ask for additional constructions:
\begin{itemize}
\item logarithmic VOAs ask for a regularized infinite free-field or
screening resolution;
\item quasi-rational VOAs ask for a canonical completed quasi-free
resolution and a conformal homotopy-invariance theorem;
\item critical-level affine algebras ask for a screening model based on
the Feigin--Frenkel centre.
\end{itemize}

\begin{conjecture}[Logarithmic conductor]
\label{conj:K-logarithmic}
\ClaimStatusConjectured
For each triplet algebra $W_p$, there exists a flat deformation chart
$A_\epsilon\to W_p$ equipped with conformal BRST presentations whose
regularized ghost characters extend to $\epsilon=0$.  The finite part
at the central fibre is invariant under conformal homotopies inside the
chart and defines $K_{\mathrm{ghost}}^{\log}(W_p)$.
\end{conjecture}

\begin{conjecture}[Quasi-rational presentation independence]
\label{conj:K-quasi-rational}
\ClaimStatusConjectured
Let $A$ be a quasi-rational VOA equipped with two completed quasi-free
conformal BRST presentations.  A filtered conformal BRST equivalence
between the presentations identifies their regularized ghost
characters and hence their scalar finite parts.
\end{conjecture}

\begin{conjecture}[Critical-level screening conductor]
\label{conj:K-critical}
\ClaimStatusConjectured
For a simple $\mathfrak g$, a completed Feigin--Frenkel screening
complex at $k=-h^\vee$ carries a regularized conformal trace whose
finite part is invariant under changes of screening generators.  This
finite part defines a critical-level scalar $K^{\mathrm{FF}}$.
\end{conjecture}

\section{Cross-volume parallel: the Borcherds scalar}
\label{sec:conductor-cross-volume}

The Vol.~III Borcherds--Kac--Moody scalar
\[
\kappa_{\mathrm{BKM}}(\mathfrak g_\Phi)=\frac{c_\Phi(0)}2
\]
is the automorphic weight of the Borcherds lift of a chosen weak Jacobi
input~$\Phi$.  Its type signature is
\[
(\text{CY quadrant},\ \text{Borcherds-product presentation},\
 \text{level }5,\ H_{\mathrm{Borch}}).
\]
For the $K3\times E$ denominator normalization,
$c_{\Delta_5}(0)=10$ and
$\kappa_{\mathrm{BKM}}(\Delta_5)=5$; the Igusa square
$\Phi_{10}^{\mathrm{un}}=\Delta_5^2$ has weight~$10$.  The same row
also carries the proved values
\[
\{\kappa_{\mathrm{cat}},\kappa_{\mathrm{ch}}^{\mathrm{Heis}},
\kappa_{\mathrm{BKM}},\kappa_{\mathrm{fiber}}\}(K3\times E)
=\{0,3,5,24\}.
\]

Write
\[
H_{\mathsf B}:=(H_{\mathrm{chart}},H_{\mathrm{KD}},H_{\mathrm{scalar}},
H_{\mathrm{mod}},H_{\mathrm{quant}}).
\]
The Mukai lattice
\[
\widetilde H(K3,\mathbb Z)
\simeq U^4\oplus E_8(-1)^2
\]
has signature $(4,20)$, so the lattice identity $2c_+=8$ is exact.
Its proposed interpretation as a chiral conductor is the
\ClaimStatusConjectured{} comparison
\[
K^{\kappa_{\mathrm{ch}}}_{\mathsf B}=2c_+=8
\qquad\text{under }H_{\mathsf B}.
\]
Bruinier~\cite[Lemma~5.1]{Bruinier2002} supplies cyclotomic
integrality for Fourier coefficients and the half-integral-weight
integer $N'=\operatorname{lcm}(N,8)$; Lusztig's root-of-unity theory
takes the root order as input.  Thus $H_{\mathrm{mod}}$ constructs the
Humbert comparison and $H_{\mathrm{quant}}$ selects order~$8$.
The Universal Trace Identity is the conjectural bridge from this
CY-quadrant package to a Vol.~I scalar Verdier conductor.
